\chapter{引言}\label{ux5f15ux8a00}

\section{研究背景与意义}\label{ux7814ux7a76ux80ccux666fux4e0eux610fux4e49}

\subsection{研究背景}\label{ux7814ux7a76ux80ccux666f}

近年来，国家层面持续推动文化和旅游深度融合。2018 年文化和旅游部组建后，``以文塑旅、以旅彰文''逐渐成为文旅融合工作的指导性表述；《``十四五''文化发展规划》《``十四五''旅游业发展规划》也对非物质文化遗产、历史文化名城名镇名村、风景名胜区与旅游景区体系的协同利用提出了要求。在这一政策背景下，区县尺度如何将文化资源家底转化为游客可感知、可参与、可体验的旅游产品，成为规划实践中需要回应的问题。

佛山市南海区位于广东省中部，是广佛同城化发展的重要区域。区内文化资源类型较为丰富：本研究整理的非遗名录共包含国家、省、市、区四级非物质文化遗产 90 项；进入核心指数计算的文化载体样本共 165 条，其中包括不可移动文物保护单位 80 处、历史文化名镇名村与传统村落 12 项、文化景观 19 项、圩市街区 18 项及非遗空间锚点 36 项。不可移动文物保护单位中，全国重点文物保护单位 2 处、省级文物保护单位 19 处、市级文物保护单位 59 处。除上述核心分析样本外，官方资源扩展底表另纳入 484 处普查口径不可移动文物，用于补充诊断与空间覆盖校验，不直接替代 165 条核心指数样本。

但是，文化资源的丰富并不必然转化为旅游体验的充分供给。在南海区现有文旅开发中，部分文化资源仍主要停留在典籍记载、保护名录或地方记忆层面，尚未稳定转化为可识别、可到达、可体验的旅游产品。文化资源与旅游活动之间的错位，在区县尺度上仍缺少可量化、可复核的识别方法，因而影响后续的资源整合、产品策划与空间治理。本研究尝试结合多源数据、大语言模型和知识图谱方法，构建面向区县文旅融合潜力识别的分析框架。

\subsection{研究意义}\label{ux7814ux7a76ux610fux4e49}

\textbf{实践意义：}研究面向南海区政府及相关规划部门，在 165 条核心文化载体样本层面给出可复核的错位诊断清单，支撑文旅资源梳理、产品策划与空间布局；并按镇街与等级提出分层级、分片区的潜力释放建议。

\textbf{方法意义：}本研究探索将大语言模型、知识图谱、多源 POI 与评论数据结合，形成一套区县级文旅融合分析流程。其核心环节包括文化记忆指数、官方认证指数、旅游热度指数与错位指数，具有一定可迁移性，可为同类尺度的区县个案提供参考，同时也可以扩展到更大区域范围内。

\section{研究目标与内容}\label{ux7814ux7a76ux76eeux6807ux4e0eux5185ux5bb9}

\subsection{研究目标}\label{ux7814ux7a76ux76eeux6807}

建立以``物质遗产为主桥梁、非遗为动态补充''的双谱系对照框架，对南海区文旅融合现状进行量化识别，并对错位的结构性来源与可干预条件作出解释，最终服务于规划层面的决策建议。

\subsection{研究内容}\label{ux7814ux7a76ux5185ux5bb9}

本研究包含六项主要工作：（1）整理典籍、POI、评论、非遗名录与政府普查文化载体等多源数据；（2）基于典籍构建文化知识图谱；（3）基于 POI、评分和评论量刻画旅游供给与热度；（4）以 165 条文化载体样本为主桥梁，构造文化记忆指数（CMI）、官方认证指数（OAI）、旅游热度指数（THI）与错位指数（MI）；（5）在 500 m 网格尺度引入知识图谱 0/1 跳路径复核，结合相关矩阵识别潜力释放条件；（6）给出分片区、分层级的规划建议。

\section{技术路线与论文结构}\label{ux6280ux672fux8defux7ebfux4e0eux8bbaux6587ux7ed3ux6784}

\subsection{技术路线}\label{ux6280ux672fux8defux7ebf}

研究按``数据采集 → 清洗与结构化 → 知识图谱构建 → 桥梁匹配 → 量化指数与错位指数 → 空间分析与潜力条件识别 → 建议输出''的链条推进。采集层包括 OCR（典籍）、API 抓取（POI 与评论）、官方 Shapefile 与名录整理；清洗层包括三源 POI 去重、坐标统一与类目纠正；建图层使用大语言模型抽取后导入 Neo4j；匹配层以 165 条文化载体样本为核心桥梁，非遗作为动态补充；量化层产出四个指数与错位分类；空间层开展核密度估计、DBSCAN 聚类与相关矩阵分析；输出层形成规划建议与附录支撑材料。

\subsection{论文结构}\label{ux8bbaux6587ux7ed3ux6784}

第 2 章为文献综述与相关研究述评；第 3 章介绍研究区域与数据基础；第 4 章说明知识图谱的构建方法与成果；第 5 章阐述耦合桥梁的主---补结构、量化指数与错位指数；第 6 章讨论空间格局、潜力释放条件、权重稳健性和评论语义补充检验；第 7 章给出结论与建议。
